\documentclass[a4paper,12pt,twoside]{memoir}

% Castellano
\usepackage[spanish,es-tabla]{babel}
\selectlanguage{spanish}
\usepackage[utf8]{inputenc}
\usepackage[T1]{fontenc}
\usepackage{lmodern} % scalable font
\usepackage{microtype}
\usepackage{placeins}

\RequirePackage{booktabs}
\RequirePackage[table]{xcolor}
\RequirePackage{xtab}
\RequirePackage{multirow}

% Links
\usepackage[colorlinks]{hyperref}
\hypersetup{
	allcolors = {red}
}

% Ecuaciones
\usepackage{amsmath}

% Rutas de fichero / paquete
\newcommand{\ruta}[1]{{\sffamily #1}}

% Párrafos
\nonzeroparskip


% Imagenes
\usepackage{graphicx}
\newcommand{\imagen}[2]{
	\begin{figure}[!h]
		\centering
		\includegraphics[width=0.9\textwidth]{#1}
		\caption{#2}\label{fig:#1}
	\end{figure}
	\FloatBarrier
}

\newcommand{\imagenflotante}[2]{
	\begin{figure}%[!h]
		\centering
		\includegraphics[width=0.9\textwidth]{#1}
		\caption{#2}\label{fig:#1}
	\end{figure}
}



% El comando \figura nos permite insertar figuras comodamente, y utilizando
% siempre el mismo formato. Los parametros son:
% 1 -> Porcentaje del ancho de página que ocupará la figura (de 0 a 1)
% 2 --> Fichero de la imagen
% 3 --> Texto a pie de imagen
% 4 --> Etiqueta (label) para referencias
% 5 --> Opciones que queramos pasarle al \includegraphics
% 6 --> Opciones de posicionamiento a pasarle a \begin{figure}
\newcommand{\figuraConPosicion}[6]{%
  \setlength{\anchoFloat}{#1\textwidth}%
  \addtolength{\anchoFloat}{-4\fboxsep}%
  \setlength{\anchoFigura}{\anchoFloat}%
  \begin{figure}[#6]
    \begin{center}%
      \Ovalbox{%
        \begin{minipage}{\anchoFloat}%
          \begin{center}%
            \includegraphics[width=\anchoFigura,#5]{#2}%
            \caption{#3}%
            \label{#4}%
          \end{center}%
        \end{minipage}
      }%
    \end{center}%
  \end{figure}%
}

%
% Comando para incluir imágenes en formato apaisado (sin marco).
\newcommand{\figuraApaisadaSinMarco}[5]{%
  \begin{figure}%
    \begin{center}%
    \includegraphics[angle=90,height=#1\textheight,#5]{#2}%
    \caption{#3}%
    \label{#4}%
    \end{center}%
  \end{figure}%
}
% Para las tablas
\newcommand{\otoprule}{\midrule [\heavyrulewidth]}
%
% Nuevo comando para tablas pequeñas (menos de una página).
\newcommand{\tablaSmall}[5]{%
 \begin{table}
  \begin{center}
   \rowcolors {2}{gray!35}{}
   \begin{tabular}{#2}
    \toprule
    #4
    \otoprule
    #5
    \bottomrule
   \end{tabular}
   \caption{#1}
   \label{tabla:#3}
  \end{center}
 \end{table}
}

%
%Para el float H de tablaSmallSinColores
\usepackage{float}

%
% Nuevo comando para tablas pequeñas (menos de una página).
\newcommand{\tablaSmallSinColores}[5]{%
 \begin{table}[H]
  \begin{center}
   \begin{tabular}{#2}
    \toprule
    #4
    \otoprule
    #5
    \bottomrule
   \end{tabular}
   \caption{#1}
   \label{tabla:#3}
  \end{center}
 \end{table}
}

\newcommand{\tablaApaisadaSmall}[5]{%
\begin{landscape}
  \begin{table}
   \begin{center}
    \rowcolors {2}{gray!35}{}
    \begin{tabular}{#2}
     \toprule
     #4
     \otoprule
     #5
     \bottomrule
    \end{tabular}
    \caption{#1}
    \label{tabla:#3}
   \end{center}
  \end{table}
\end{landscape}
}

%
% Nuevo comando para tablas grandes con cabecera y filas alternas coloreadas en gris.
\newcommand{\tabla}[6]{%
  \begin{center}
    \tablefirsthead{
      \toprule
      #5
      \otoprule
    }
    \tablehead{
      \multicolumn{#3}{l}{\small\sl continúa desde la página anterior}\\
      \toprule
      #5
      \otoprule
    }
    \tabletail{
      \hline
      \multicolumn{#3}{r}{\small\sl continúa en la página siguiente}\\
    }
    \tablelasttail{
      \hline
    }
    \bottomcaption{#1}
    \rowcolors {2}{gray!35}{}
    \begin{xtabular}{#2}
      #6
      \bottomrule
    \end{xtabular}
    \label{tabla:#4}
  \end{center}
}

%
% Nuevo comando para tablas grandes con cabecera.
\newcommand{\tablaSinColores}[6]{%
  \begin{center}
    \tablefirsthead{
      \toprule
      #5
      \otoprule
    }
    \tablehead{
      \multicolumn{#3}{l}{\small\sl continúa desde la página anterior}\\
      \toprule
      #5
      \otoprule
    }
    \tabletail{
      \hline
      \multicolumn{#3}{r}{\small\sl continúa en la página siguiente}\\
    }
    \tablelasttail{
      \hline
    }
    \bottomcaption{#1}
    \begin{xtabular}{#2}
      #6
      \bottomrule
    \end{xtabular}
    \label{tabla:#4}
  \end{center}
}

%
% Nuevo comando para tablas grandes sin cabecera.
\newcommand{\tablaSinCabecera}[5]{%
  \begin{center}
    \tablefirsthead{
      \toprule
    }
    \tablehead{
      \multicolumn{#3}{l}{\small\sl continúa desde la página anterior}\\
      \hline
    }
    \tabletail{
      \hline
      \multicolumn{#3}{r}{\small\sl continúa en la página siguiente}\\
    }
    \tablelasttail{
      \hline
    }
    \bottomcaption{#1}
  \begin{xtabular}{#2}
    #5
   \bottomrule
  \end{xtabular}
  \label{tabla:#4}
  \end{center}
}



\definecolor{cgoLight}{HTML}{EEEEEE}
\definecolor{cgoExtralight}{HTML}{FFFFFF}

%
% Nuevo comando para tablas grandes sin cabecera.
\newcommand{\tablaSinCabeceraConBandas}[5]{%
  \begin{center}
    \tablefirsthead{
      \toprule
    }
    \tablehead{
      \multicolumn{#3}{l}{\small\sl continúa desde la página anterior}\\
      \hline
    }
    \tabletail{
      \hline
      \multicolumn{#3}{r}{\small\sl continúa en la página siguiente}\\
    }
    \tablelasttail{
      \hline
    }
    \bottomcaption{#1}
    \rowcolors[]{1}{cgoExtralight}{cgoLight}

  \begin{xtabular}{#2}
    #5
   \bottomrule
  \end{xtabular}
  \label{tabla:#4}
  \end{center}
}




\graphicspath{ {./img/} }

% Capítulos
\chapterstyle{bianchi}
\newcommand{\capitulo}[2]{
	\setcounter{chapter}{#1}
	\setcounter{section}{0}
	\chapter*{#2}
	\addcontentsline{toc}{chapter}{#2}
	\markboth{#2}{#2}
}

% Apéndices
\renewcommand{\appendixname}{Apéndice}
\renewcommand*\cftappendixname{\appendixname}

\newcommand{\apendice}[1]{
	%\renewcommand{\thechapter}{A}
	\chapter{#1}
}

\renewcommand*\cftappendixname{\appendixname\ }

% Formato de portada
\makeatletter
\usepackage{xcolor}
\newcommand{\tutor}[1]{\def\@tutor{#1}}
\newcommand{\course}[1]{\def\@course{#1}}
\definecolor{cpardoBox}{HTML}{E6E6FF}
\def\maketitle{
  \null
  \thispagestyle{empty}
  % Cabecera ----------------
\noindent\includegraphics[width=\textwidth]{cabecera}\vspace{1cm}%
  \vfill
  % Título proyecto y escudo informática ----------------
  \colorbox{cpardoBox}{%
    \begin{minipage}{.8\textwidth}
      \vspace{.5cm}\Large
      \begin{center}
      \textbf{TFG del Grado en Ingeniería Informática}\vspace{.6cm}\\
      \textbf{\LARGE\@title{}}
      \end{center}
      \vspace{.2cm}
    \end{minipage}

  }%
  \hfill\begin{minipage}{.20\textwidth}
    \includegraphics[width=\textwidth]{escudoInfor}
  \end{minipage}
  \vfill
  % Datos de alumno, curso y tutores ------------------
  \begin{center}%
  {%
    \noindent\LARGE
    Presentado por \@author{}\\ 
    en Universidad de Burgos --- \@date{}\\
    Tutor: \@tutor{}\\
  }%
  \end{center}%
  \null
  \cleardoublepage
  }
\makeatother


% Datos de portada
\title{título del TFG \\Documentación Técnica}
\author{nombre alumno}
\tutor{nombre tutor}
\date{\today}

\begin{document}

\maketitle



\cleardoublepage



%%%%%%%%%%%%%%%%%%%%%%%%%%%%%%%%%%%%%%%%%%%%%%%%%%%%%%%%%%%%%%%%%%%%%%%%%%%%%%%%%%%%%%%%



\frontmatter


\clearpage

% Indices
\tableofcontents

\clearpage

\listoffigures

\clearpage

\listoftables

\clearpage

\mainmatter

\appendix

\apendice{Plan de Proyecto Software}

\section{Introducción}
Para que un proyecto se desarrolle con normalidad y con el menor número de imprevistos posibles es esencial que cuente con una fase de planificación. Aquí se estima el tiempo, trabajo y dinero que puede llegar a usarse durante la realización del proyecto. Para ello, se debe analizar en detalle cada parte del proyecto. De cara al futuro, el análisis del proyecto puede servir para predecir como de bien puede desarrollarse una continuación del mismo.

La planificación del proyecto consta de dos partes:
\begin{itemize}
	\tightlist
	\item \textbf{Planificación temporal}: en esta parte se analiza y planifica el tiempo que se va a dedicar a cada parte del proyecto, fecha de inicio y final aproximado, teniendo en cuenta el trabajo necesario para cada parte.
	\item \textbf{Estudio de viabilidad}: en esta parte se analiza como de viable es la realización del proyecto, se divide a su vez en dos apartados:
	
	\begin{itemize}
		\tightlist
		\item \textbf{Económica}: en esta parte se estiman los costes y los beneficios que puede suponer el proyecto.
		\item \textbf{Legal}: en esta parte se analizan los conceptos legales del proyecto, como podrían ser las licencias del proyecto o la política de protección de datos.
	\end{itemize}
\end{itemize}

\section{Planificación temporal}
La planificación temporal se organiza mediante \textit{sprints}. Cada \textit{sprint} dura una o dos semanas. Los \textit{sprints} empiezan en la reunión y acaban en la siguiente.\\

\subsection{Sprints 0-2}
Estos \textit{sprints} se centran en investigar las tecnologías que se van a usar en el proyecto, toma de contacto con las herramientas, construcción de prototipos en las tecnologías y desplegar la aplicación básica.\\

Se comentó con los tutores que en proyectos webs anteriores el despliegue se solía dejar para las etapas finales del proyecto, haciendo que todos los problemas asociados surjan en esas finales, retrasando el despliegue y, a veces, no llegar a desplegar la aplicación. Por eso se ha desplegado la aplicación desde las etapas iniciales del proyecto.\\

\subsection{Sprint 3}
Este \textit{sprint} se centra en mejorar visualmente la aplicación y encabezar en desarrollo del proyecto hacia algo definitivo. Se mejora la estructura lógica del proyecto, la visualización de un espectro individual, se añade la subida y visualización de \textit{datasets} y soporte de usuarios mediante cuenta de Google.\\

Hubo dos tareas que no se llegaron a completar, la subida y visualización y el manual de despliegue definitivo.\\

\subsection{Sprint 4}
En este \textit{sprint} se completan las tareas que no habían dado tiempo del \textit{sprint} anterior y se planteó usar una base de datos en lugar de almacenamiento para guardar los datos. También se añadió la opción de borrar un \textit{dataset} ya almacenado.\\

En este \textit{sprint} se han completado todas las tareas.\\

\subsection{Sprint 5}
De la comparación del \textit{sprint} anterior se decidió usar MongoDB para el almacenamiento, por lo cual la mayoría de los esfuerzos se centraron en adaptar la aplicación para usar MongoDB en todos sus aspectos: guardar, borrar y coger los datos.\\

En este \textit{sprint} se han completado todas las tareas.\\

\subsection{Sprint 6}
Para este \textit{sprint} se planteó cambiar la forma en la que se guardan los \textit{datasets} para que más adelante sea más sencillo de pasar los datos al sistema de aprendizaje automático (muchos dataframes a un dataframe). También se añadió soporte de comentarios en el \textit{dataset} y se añadió soporte para un fichero excel en la subida que contenga los metadatos del \textit{dataset}. Por último se empezó a añadir controles en la parte de visualización para el procesado de los espectros.\\

Esta última tarea no se pudo completar.\\

\subsection{Sprint 7}
Este \textit{sprint} se centró en completar la tarea del \textit{sprint} anterior, mostrar en la parte de visualización una tabla con los ejemplos del \textit{dataset} y sus metadatos, documentar el código y avanzar en la memoria y anexos.\\

\section{Estudio de viabilidad}

\subsection{Viabilidad económica}

\subsection{Viabilidad legal}



\apendice{Especificación de Requisitos}

\section{Introducción}

En este apéndice se describen los objetivos generales de la aplicación y se 
detallan sus requisitos, tanto funcionales como no funcionales.

\section{Objetivos generales}

\begin{itemize}
	\item Ofrecer control de usuarios.
	\item Permitir a los usuarios subir \textit{datasets} y espectros.
	\item Que los \textit{datasets} y espectros se puedan visualizar.
	\item Que sobre la visualización se pueda aplicar operaciones de 
	procesamiento.
	\item Poder entrenar modelos de aprendizaje automático con los
	\textit{datasets} subidos.
	\item Poder usar los modelos mencionados anteriormente para predecir nuevos
	ejemplos subidos.
	\item Que la aplicación final sea útil para la investigadora, que hace las 
	veces de cliente en este proyecto.
\end{itemize}

\section{Catálogo de requisitos}

\subsection{Requisitos funcionales}

\begin{itemize}
	\item \textbf{RF-1 Control de usuarios}: la aplicación debe permitir 
	controlar usuarios.
	\begin{itemize}
		\item \textbf{RF-1.1 Integración con Google}: la aplicación debe poder 
		hacer uso de cuentas de Google para el control de usuarios.
		\item \textbf{RF-1.2 Inicio de sesión}: el usuario debe poder iniciar 
		sesión con una	cuenta de Google.
		\item \textbf{RF-1.3 Cierre de sesión}: el usuario debe poder cerrar 
		sesión cuando haya terminado.
	\end{itemize}
	\item \textbf{RF-2 Datasets}: la aplicación debe poder almacenar y 
	gestionar conjuntos de espectros.
	\begin{itemize}
		\item \textbf{RF-2.1 Subida}: el usuario debe poder subir un 
		\textit{dataset}.
		\item \textbf{RF-2.2 Eliminado}: el usuario debe poder eliminar un 
		\textit{dataset} almacenado.
		\item \textbf{RF-2.3 Visualización}: el usuario debe poder visualizar 
		un \textit{dataset} almacenado.
	\end{itemize}
	\item \textbf{RF-3 Espectros}: la aplicación debe poder almacenar y 
	gestionar espectros.
	\begin{itemize}
		\item \textbf{RF-3.1 Subida}: el usuario debe poder subir un espectro.
		\item \textbf{RF-3.2 Eliminado}: el usuario debe poder eliminar un 
		espectro.
		\item \textbf{RF-3.3 Visualización}: el usuario debe poder visualizar 
		un espectro.
	\end{itemize}
	\item \textbf{RF-4 Procesamiento}: el usuario debe poder aplicar 
	operaciones de preprocesamiento.
	\begin{itemize}
		\item \textbf{RF-4.1: Procesamiento de \textit{dataset}}: el usuario 
		debe poder aplicar operaciones de preprocesamiento sobre un 
		\textit{dataset} visualizado.
		\item \textbf{RF-4.2: Procesamiento de espectro}: el usuario debe poder 
		aplicar operaciones de preprocesamiento sobre un espectro visualizado.
	\end{itemize}
	\item \textbf{RF-5 Minería de datos}: el usuario debe poder usar técnicas 
	de minería de datos.
	\begin{itemize}
		\item \textbf{RF-5.1 Creación}: el usuario debe poder crear modelos 
		personalizados.
		\item \textbf{RF-5.2 Evaluación}: la aplicación debe ofrecer métricas 
		del modelo entrenado.
		\item \textbf{RF-5.3 Predicción}: el usuario debe poder usar los 
		modelos que ha creado para predecir nuevos espectros.
	\end{itemize}
\end{itemize}

\subsection{Requisitos no funcionales}

	\begin{itemize}
	\item \textbf{RNF-1 Usabilidad}: la aplicación debe ser intuitiva y fácil 
	de usar.
	\item \textbf{RNF-2 Escalabilidad}: el rendimiento debe poder aumentar al 
	incrementar los recursos.
	\item \textbf{RNF-3 Manteninibilidad}: debe ser sencillo añadir 
	funcionalidad nueva a la aplicación.
	\item \textbf{RNF-4 Compatibilidad}: la aplicación debe poder funcionar en 
	los principales navegadores.
	\item \textbf{RNF-5 Responsividad}: la aplicación debe adaptarse al tamaño 
	de la pantalla.
	\item \textbf{RNF-6 Facilidad de despliegue}: la aplicación debe poder 
	desplegarse en un servidor de forma sencilla.
	\end{itemize}

\section{Especificación de requisitos}

En esta sección se presentan el diagrama de casos de uso. A continuación, se 
desarrollan los casos de uso mostrados en el diagrama.

\imagen{casos-uso}{Diagrama de casos de uso}

\casoDeUso{1}{Iniciar sesión}
{RF-1, RF-1.1, RF-1.2}
{El usuario inicia sesión en la aplicación.}
{Navegador web abierto y página de la aplicación cargada.}
{
	\item Pulsar botón ``Iniciar sesión con Google''.
	\item Introducir o seleccionar cuenta de Google.
}
{Redirección a la aplicación con sesión iniciada.}
{	\item Cuenta no existente.
	\item Combinación de nombre y contraseña incorrecta.
}
{Alta}

\casoDeUso{2}{Cerrar sesión}
{RF-1, RF-1.3}
{El usuario cierra sesión en la aplicación.}
{Sesión iniciada en la aplicación.}
{
	\item Pulsar en el botón cuyo texto es el correo.
	\item Pulsar ``Cerrar sesión''.
}
{Redirección a la aplicación con sesión cerrada.}
{	\item La sesión no estaba iniciada.
}
{Alta}

\casoDeUso{3}{Subir dataset}
{RF-2, RF-2.1}
{El usuario sube un dataset a la aplicación para su guardado.}
{Sesión iniciada en la aplicación.}
{
	\item Pulsar en el botón ``Subir dataset''.
	\item Descargar y rellenar la plantilla.
	\item Crear un fichero \code{.zip} con los datos y la plantilla.
	\item Rellenar el formulario de subida.
	\item Seleccionar el fichero creado.
	\item Presionar el botón ``Subir''.
}
{Redirección a la página de los archivos guardados.}
{	\item El formato del dataset no es correcto.
	\item Existe un dataset con el mismo nombre.
}
{Alta}

\casoDeUso{4}{Subir espectro}
{RF-3, RF-3.1}
{El usuario sube un espectro a la aplicación para su guardado.}
{Sesión iniciada en la aplicación.}
{
	\item Pulsar en el botón ``Subir espectro''.
	\item Rellenar el formulario de subida.
	\item Seleccionar el espectro.
	\item Presionar botón ``Subir''.
}
{Redirección a la página de los archivos guardados.}
{	\item El formato del espectro no es correcto.
	\item Existe un espectro con el mismo nombre.
}
{Alta}

\casoDeUso{5}{Visualizar dataset}
{RF-2, RF-2.3, RF-4, RF-4.1}
{El usuario visualiza un espectro en la aplicación.}
{Sesión iniciada en la aplicación, dataset guardado en la aplicación y estar en 
la página de ``Mis ficheros''.}
{
	\item Escoger un dataset que visualizar.
	\item Pulsar el botón ``Visualizar'' en el dataset escogido.
}
{Se muestra la página de visualización.}
{	\item No se ha podido cargar el dataset.
}
{Alta}

\casoDeUso{6}{Visualizar espectro}
{RF-3, RF-3.3, RF-4, RF-4.2}
{El usuario visualiza un espectro en la aplicación.}
{Sesión iniciada en la aplicación, espectro guardado en la aplicación y estar 
en la página de ``Mis ficheros''.}
{
	\item Escoger un espectro que visualizar.
	\item Pulsar el botón ``Visualizar'' en el espectro escogido.
}
{Se muestra la página de visualización.}
{	\item No se ha podido cargar el espectro.
}
{Alta}

\casoDeUso{7}{Guardar clasificador}
{RF-5, RF-5.1, RF-5.2}
{El usuario crear y entrena un modelo usando como datos un dataset guardado.}
{Sesión iniciada en la aplicación, dataset guardado en la aplicación y estar en 
la página de ``Mis ficheros''.}
{
	\item Escoger el dataset que se quiera usar como base.
	\item Pulsar el botón ``Crear modelo'' en el dataset escogido.
	\item Seleccionar en el desplegable el modelo a usar.
	\item (Opcional) Rellenar el formulario con los parámetros del modelo.
	\item Pulsar el botón ``Crear y evaluar modelo''.
	\item Para guardar el clasificador:
	\begin{enumerate}
		\item Completar el formulario.
		\item Pulsar el botón ``Guardar''.
	\end{enumerate}
	\item Para no guardar el clasificador:
	\begin{enumerate}
		\item Pulsar el botón ``Descartar''.
	\end{enumerate}
}
{Se redirige a los archivos guardados.}
{	\item Los parámetros del modelo introducidos son erróneos.
	\item Ya existe un clasificador con el nombre introducido.
}
{Media}

\casoDeUso{8}{Predecir espectro}
{RF-5, RF-5.3}
{El usuario usa un clasificador creado para predecir un espectro.}
{Sesión iniciada en la aplicación, dataset guardado en la aplicación, algún 
clasificador creado y estar en la página de ``Mis ficheros''.}
{
	\item Escoger el espectro que se quiera predecir.
	\item Pulsar el botón ``Predecir'' en el espectro escogido.
	\item Seleccionar en el desplegable el clasificador a usar.
	\item Pulsar el botón ``Predecir espectro''.
	\item Para guardar el clasificador:
}
{Se muestra la predicción.}
{	\item No se puede cargar el clasificador.
}
{Media}


\apendice{Especificación de diseño}

\section{Introducción}

En este apéndice se expone el diseño que ha dado lugar a la aplicación. Se 
incluye el diseño de datos, diseño procedimental y diseño arquitectónico.

\section{Diseño de datos}

En esta sección se explican las entidades usadas por la aplicación:

\begin{itemize}
	\item Dataset: esta entidad representa los conjuntos de espectros subidos 
	por los usuarios.
	\item Spectrum: esta entidad representa un espectro subido por los usuarios.
	\item ClassifierSet: esta entidad representa los clasificadores entrenados 
	por los usuarios, se compone de tres clasificadores, uno por cada atributo 
	a predecir.
\end{itemize}

\imagen{clases}{Diagrama de clases}

\section{Diseño procedimental}

En esta sección se explican las interacciones más importantes de la aplicación. 
Las interacción principal de la aplicación es, dentro de la visualización de 
datasets, cuando se selecciona un espectro para visualizar, ya que entran en 
juego todas las partes de la aplicación (ver figura~\ref{fig:carga-espectro}).

\imagen{carga-espectro}{Diagrama de secuencia de visualización de espectro}

Cuando el usuario selecciona un espectro, se envía al servidor la fila 
seleccionada, el servidor pide a la base de datos el dataset que se está 
visualizando y lo devuelve, el servidor extrae del dataset el espectro 
seleccionado, lo convierte en una figura que se pueda visualizar y lo envía a 
la interfaz, que se lo muestra al usuario.

A continuación, para mostrar el espectro procesado, la interfaz envía al 
servidor los valores actuales de los controles, el servidor aplica las acciones 
de procesamiento al espectro, lo convierte en una figura para poder 
visualizarlo y lo envía a la interfaz, que lo muestra al usuario.

\section{Diseño arquitectónico}

Al haber desarrollado una aplicación web, la arquitectura del proyecto está 
condicionada por ello.

\subsection{Modelo Vista Controlador (MVC)}

Se ha seguido el patrón de MVC con el objetivo de separar la lógica de la 
aplicación (controlador), los datos (modelo) y la interfaz (vista). Este patrón 
permite facilitar la tarea del mantenimiento. Además, al ser un patrón 
conocido, se facilita la tarea de trabajar en la aplicación a futuros 
desarrolladores. 

\imagen{mvc}{Diagrama del patrón MVC~\cite{wiki:mvc}}

Los modelos se encuentran presentes en \code{SpectraViewer/models.py}, los 
controladores se encuentran en \code{SpectraViewer/main/routes.py} y las vistas 
se encuentran en \code{SpectraViewer/templates/}.

\apendice{Manual del programador}

\section{Introducción}

\section{Estructura de directorios}

\section{Manual del programador}

\subsection{Manual de despliegue}
El despliegue de la aplicación se hace mediante la herramienta Nanobox, cuyo
propósito principal es el de facilitar la tarea de despliegue. Para ello combina
servicios de virtualización con servidores virtuales privados o VPS, ocupándose
de la configuración de la máquina virtual que el proveedor de VPS nos
proporcione.\\

El primer paso es registrarse en alguno de los proveedores
disponibles\footnote{\url{https://docs.nanobox.io/providers/hosting-accounts/}}.
Para este proyecto se ha elegido Digital Ocean debido a que durante el proyecto
se contaba con el ``Student Developer Pack'' de GitHub, el cual contiene un
crédito gratuito de 50\$ para esa plataforma.\\

El siguiente paso es crear una cuenta en
Nanobox\footnote{\url{https://nanobox.io/}}. Una vez hecho, hay que enlazar esta
cuenta con la creada anteriormente, desde las opciones de la cuenta en la
pestaña ``Hosting Accounts'', una vez ahí seguir las instrucciones que se
muestran en la página.\\

Después de tener las cuentas preparadas hay que descargar la herramienta Nanobox
para el sistema adecuado e instalarla, seguir las instrucciones de la
documentación oficial\footnote{\url{https://docs.nanobox.io/install/}} y las del
instalador.\\

Desde la página principal se pulsa en ``Lauch New App'' para crear la nueva
aplicación, seguir las instrucciones en la página. Al terminar se muestran
instrucciones para desplegar la aplicación.\\

Como se puede comprobar, Nanobox cumple el propósito de facilitar el despliegue.
Para cualquier otra duda sobre la herramienta toda la documentación oficial está
disponible en \url{https://docs.nanobox.io/}.\\
\section{Compilación, instalación y ejecución del proyecto}

\section{Pruebas del sistema}


\apendice{Documentación de usuario}

\section{Introducción}

En este apéndice se explica los requisitos que debe cumplir el usuario para 
ejecutar la aplicación, como lanzarla y como usarla.

\section{Requisitos de usuarios}

Al tratarse de una aplicación los requisitos que debe cumplir el usuario son 
los siguientes:
\begin{itemize}
	\tightlist
	\item Navegador web instalado.
	\item Cuenta de Google activa.
	\item JavaScript avtivo en el navegador.
	\item \textit{Cookies} activas en el navegador.
\end{itemize}

La aplicación se ha probado en los siguientes navegadores y se certifica que 
funciona:
\begin{itemize}
	\tightlist
	\item Google Chrome 67.0.3396.99.
	\item Chromium 66.0.3359.181.
	\item Mozilla Firefox 60.0.2.
\end{itemize}

Aunque la aplicación funciona en dispositivos móvil no se recomienda hacerlo 
debido a que la visualización sería demasiado pequeña y no estarían disponibles 
las acciones sobre los gráficos.

\section{Instalación}

Debido a que se proporciona una aplicación web no es necesario instalarla para 
poder usarla. Sin embargo, si se quiere proceder a la instalación se pueden 
seguir las instrucciones en la sección~\ref{sec:instalacion}.

De cara a probar la aplicación con ejemplos ya cargados se proporciona un 
cuenta con ejemplos cargados para su uso. El usuario proporcionado es  
\code{tfg.visor.ejemplos@gmail.com} y su contraseña es 
\code{tfg\_visor\_ejemplos}.

Hay que tener en cuenta que al ser una cuenta de prueba la pueden usar varias 
personas por lo que se recomienda encarecidamente no borrar los espectros ni 
los datasets que se encuentran subidos.

\section{Manual del usuario}

En esta sección se enseña al usuario como manejar la aplicación.

\subsection{Inicio}

Nada más entrar a la aplicación se muestra la bienvenida. Si no se ha iniciado 
sesión aparece el mensaje de la imagen~\ref{fig:manual-uso/bienvenida}.

\imagen{manual-uso/bienvenida}{Bienvenida}

Después de iniciar sesión (botón en la esquina superior derecha), o si no se 
había cerrado sesión anteriormente, se muestra el mensaje de la 
imagen~\ref{fig:manual-uso/bienvenida-inicio}.

\imagen{manual-uso/bienvenida-inicio}{Bienvenida después de iniciar sesión}

Las opciones presentadas en la barra de navegación y en las tarjetas situadas 
debajo del mensaje son las mismas, con la excepción de que en las tarjetas 
aparece una pequeña descripción de la acción.

\subsection{Mis ficheros}

Pulsando en ``Mis ficheros'' nos aparece la página con todo lo que tenemos 
asociado en nuestra cuenta (ver figura~\ref{fig:manual-uso/ficheros}), 
organizado en tres columnas, datasets, espectros y clasificadores. Para cada 
entidad se muestra el nombre, los comentarios y las opciones disponibles.

\imagen{manual-uso/ficheros}{Mis ficheros}

\subsection{Subir dataset}

Pulsando en ``Subir dataset'' nos dirigimos a la página en la que podemos subir 
en conjunto de espectros. A la izquierda se muestran las instrucciones a seguir 
y unas notas respecto al formulario, presente a la derecha. En la 
imagen~\ref{fig:manual-uso/subir-dataset} se puede ver la página con el 
formulario completo.

\imagen{manual-uso/subir-dataset}{Página para subir un dataset}

 Después de pulsar el botón ``Subir'' se muestra un mensaje 
de espera mientra se procesa la petición (ver 
figura~\ref{fig:manual-uso/espera}).

\imagen{manual-uso/espera-dataset}{Mensaje de espera al subir un dataset}

Cuando la subida del fichero termine, mostrará la página con los ficheros 
guardados, indicando que el dataset se ha subido correctamente.

\subsection{Subir espectro}

Pulsando en ``Subir espectro'' nos dirigimos a la página en la que podemos 
subir un espectro. A la izquierda se muestran las indicaciones respecto al 
formato requerido en el fichero del espectros. En la 
imagen~\ref{fig:manual-uso/subir-espectro} se puede ver la página con el 
formulario completo.

\imagen{manual-uso/subir-espectro}{Página para subir un espectro}

Después de pulsar el botón ``Subir'' se muestra un mensaje 
de espera mientra se procesa la petición (ver 
figura~\ref{fig:manual-uso/espera}). Sin embargo, como esta operación suele 
tardar poco tiempo, el mensaje no llega a leerse.

\imagen{manual-uso/espera-espectro}{Mensaje de espera al subir un espectro}

Cuando la subida del fichero termine, mostrará la página con los ficheros 
guardados, indicando que el espectro se ha subido correctamente.

\subsection{Eliminado}

Cuando se quiera eliminar cualquiera de los elementos guardados, es suficiente 
con presionar el botón ``Eliminar''. A continuación se muestra una ventana para 
confirmar la acción. Si se acepta la eliminación, se borra el elemento de la 
cuenta y se muestra un mensaje indicando que la acción se ha completado.

\subsection{Visualización y procesamiento}

Para visualizar un dataset o espectro hay que pulsar en el botón ``Visualizar'' 
en el elemento que se desee. Esta acción nos redirige a la vista de 
visualización.

\subsubsection{Datasets}

La visualización de datasets se compone de cuatro elementos 
(figura~\ref{fig:manual-uso/vis-dataset}), la tabla con los espectros 
contenidos, los controles de procesamiento, la visualización del 
espectro original y la visualización del espectro procesado. Se proporciona una 
ayuda integrada en esta vista.

\imagen{manual-uso/vis-dataset}{Visualización de dataset}

Para visualizar un espectro hay que seleccionarlo en la tabla. Pasados unos 
segundos, el espectro original y procesado aparecen en sus gráficos 
correspondientes. La aplicación soporta la visualización de varios espectros a 
la vez para poder compararlos (ver 
figura~\ref{fig:manual-uso/dataset-visualizado}).

\imagen{manual-uso/dataset-visualizado}{Visualización de dataset con espectros}

El espectro procesado se actualiza automáticamente al modificar las opciones de 
procesamiento.

\subsubsection{Espectro}

La visualización de espectro se compone de tres elementos (ver figura~\ref{f})
(figura~\ref{fig:manual-uso/vis-espectro}), los controles de procesamiento, la 
visualización del espectro original y la visualización del espectro procesado. 
Se proporciona una ayuda integrada en esta vista.

\imagen{manual-uso/vis-espectro}{Visualización de espectro}

El espectro procesado se actualiza automáticamente al modificar las opciones de 
procesamiento.

\subsection{Creación de modelos}

Para crear un modelo primero hay que escoger que dataset se quiere usar como 
referencia. Una vez decidido, se pulsa en el botón ``Crear modelo'', esto 
redirige a la página de creación de modelos (ver 
figura~\ref{fig:manual-uso/seleccionar-modelo}). 

\imagen{manual-uso/seleccionar-modelo}{Página de creación de modelos}

Para poder crear el modelo hay que seleccionar uno de los disponibles en el 
desplegable, esta acción hace visible un formulario con los posibles parámetros 
del modelo (ver figura~\ref{fig:manual-uso/modelo-seleccionado} El 
formulario se actualiza cada vez que se cambia el modelo seleccionado.

\imagen{manual-uso/modelo-seleccionado}{Formulario de creación de 
modelos}

Cuando se haya terminado de rellenar el formulario, todos los campos son 
opcionales, hay que pulsar el botón ``Crear y evaluar modelo''. Esto crea el 
modelo con los parámetros introducidos, lo entrena con el dataset seleccionado 
y lo evalúa. Mientras se realiza este proceso, se le muestra un mensaje de 
espera al usuario hasta que se completa la acción (ver 
figura~\ref{fig:manual-uso/espera-modelo}).

\imagen{manual-uso/espera-modelo}{Mensaje de espera en la creación de modelos}

Cuando se termina el entrenamiento, se muestra una página con los resultados de 
la evaluación y un formulario en caso de querer guardar el clasificador (ver 
figura~\ref{fig:manual-uso/resultados}). Ambos botones, ``Guardar'' y 
``Descartar'', llevan de vuelta a la página con los ficheros guardados, pero al 
pulsar el de guardar, guarda el clasificador en la cuenta del usuario mientras 
que pulsar el botón de descarte no.

\imagen{manual-uso/resultados}{Resultados de la evaluación}




\bibliographystyle{plain}
\bibliography{bibliografiaAnexos}

\end{document}
